\documentclass[10pt,parskip=half-]{scrartcl}
% Shared preamble for the Anspruch dossier and its companion documents.
% Load after \documentclass with % Shared preamble for the Anspruch dossier and its companion documents.
% Load after \documentclass with % Shared preamble for the Anspruch dossier and its companion documents.
% Load after \documentclass with \input{preamble} (from docs/dossier/)
% or \input{../preamble} (from docs/dossier/antrag/ or docs/dossier/partner/).
% Each document sets its own \documentclass and its own pdftitle via
% \hypersetup{pdftitle={...}}. Only packages that ship with MiKTeX.

\usepackage[T1]{fontenc}
\usepackage[utf8]{inputenc}
\usepackage{lmodern}
\usepackage[english,ngerman]{babel} % last option = main language: ngerman
\usepackage[a4paper,margin=25mm]{geometry}
\usepackage{booktabs}
\usepackage{longtable}
\usepackage{array}
\usepackage{enumitem}
\usepackage{xcolor}
\usepackage{microtype}
\usepackage[autostyle=true,german=quotes]{csquotes}
\usepackage{marginnote}
\usepackage[hidelinks,unicode,pdflang={de-DE}]{hyperref}

% ---------------------------------------------------------------------------
% Colours
\definecolor{blankbg}{RGB}{255,236,153}   % highlighted placeholder
\definecolor{todofg}{RGB}{170,40,40}      % "to confirm" margin note

% \blank{Text}: a highlighted placeholder "[Text]" for values the sender
% fills in before sending. Usable in running text and in table cells.
\newcommand{\blank}[1]{\begingroup\setlength{\fboxsep}{1pt}%
  \colorbox{blankbg}{\strut[#1]}\endgroup}

% \todo{Text}: a "zu bestätigen" note in the margin. Works in running text,
% lists and longtable cells (via marginnote, not \marginpar).
% Defined with \providecommand so that it does not clash if a document
% loads todonotes on its own (todonotes then wins).
\providecommand{\todo}[1]{\marginnote{\raggedright\scriptsize
  \textcolor{todofg}{\textbf{Zu bestätigen:} #1}}}
\setlength{\marginparwidth}{20mm}

% \draftmark: header text for every page of a draft.
\newcommand{\draftmark}{Entwurf, 23.10.2026}

% Draft header for scrartcl / scrreprt (KOMA) and article / report.
% Letters (scrlttr2, letter) place \draftmark themselves.
\makeatletter
\@ifclassloaded{scrartcl}{\def\anspruch@koma{}}{}
\@ifclassloaded{scrreprt}{\def\anspruch@koma{}}{}
\@ifclassloaded{article}{\def\anspruch@std{}}{}
\@ifclassloaded{report}{\def\anspruch@std{}}{}
\ifdefined\anspruch@koma
  \usepackage[headsepline]{scrlayer-scrpage}
  \clearpairofpagestyles
  \ohead{\small\draftmark}
  \ihead{\small\headmark}
  \cfoot*{\pagemark}
  \automark[section]{section}
  \pagestyle{scrheadings}
\fi
\ifdefined\anspruch@std
  \usepackage{fancyhdr}
  \pagestyle{fancy}
  \fancyhf{}
  \fancyhead[R]{\small\draftmark}
  \fancyhead[L]{\small\nouppercase{\leftmark}}
  \fancyfoot[C]{\thepage}
  \fancypagestyle{plain}{\fancyhf{}\fancyhead[R]{\small\draftmark}%
    \fancyfoot[C]{\thepage}\renewcommand{\headrulewidth}{0pt}}
\fi
\makeatother

% Lists: compact, as in a German Antrag.
\setlist{itemsep=0.2em,topsep=0.4em}

% Tables: a bit more air, ragged-right paragraph columns.
\renewcommand{\arraystretch}{1.15}
\newcolumntype{L}[1]{>{\raggedright\arraybackslash}p{#1}}
 (from docs/dossier/)
% or \input{../preamble} (from docs/dossier/antrag/ or docs/dossier/partner/).
% Each document sets its own \documentclass and its own pdftitle via
% \hypersetup{pdftitle={...}}. Only packages that ship with MiKTeX.

\usepackage[T1]{fontenc}
\usepackage[utf8]{inputenc}
\usepackage{lmodern}
\usepackage[english,ngerman]{babel} % last option = main language: ngerman
\usepackage[a4paper,margin=25mm]{geometry}
\usepackage{booktabs}
\usepackage{longtable}
\usepackage{array}
\usepackage{enumitem}
\usepackage{xcolor}
\usepackage{microtype}
\usepackage[autostyle=true,german=quotes]{csquotes}
\usepackage{marginnote}
\usepackage[hidelinks,unicode,pdflang={de-DE}]{hyperref}

% ---------------------------------------------------------------------------
% Colours
\definecolor{blankbg}{RGB}{255,236,153}   % highlighted placeholder
\definecolor{todofg}{RGB}{170,40,40}      % "to confirm" margin note

% \blank{Text}: a highlighted placeholder "[Text]" for values the sender
% fills in before sending. Usable in running text and in table cells.
\newcommand{\blank}[1]{\begingroup\setlength{\fboxsep}{1pt}%
  \colorbox{blankbg}{\strut[#1]}\endgroup}

% \todo{Text}: a "zu bestätigen" note in the margin. Works in running text,
% lists and longtable cells (via marginnote, not \marginpar).
% Defined with \providecommand so that it does not clash if a document
% loads todonotes on its own (todonotes then wins).
\providecommand{\todo}[1]{\marginnote{\raggedright\scriptsize
  \textcolor{todofg}{\textbf{Zu bestätigen:} #1}}}
\setlength{\marginparwidth}{20mm}

% \draftmark: header text for every page of a draft.
\newcommand{\draftmark}{Entwurf, 23.10.2026}

% Draft header for scrartcl / scrreprt (KOMA) and article / report.
% Letters (scrlttr2, letter) place \draftmark themselves.
\makeatletter
\@ifclassloaded{scrartcl}{\def\anspruch@koma{}}{}
\@ifclassloaded{scrreprt}{\def\anspruch@koma{}}{}
\@ifclassloaded{article}{\def\anspruch@std{}}{}
\@ifclassloaded{report}{\def\anspruch@std{}}{}
\ifdefined\anspruch@koma
  \usepackage[headsepline]{scrlayer-scrpage}
  \clearpairofpagestyles
  \ohead{\small\draftmark}
  \ihead{\small\headmark}
  \cfoot*{\pagemark}
  \automark[section]{section}
  \pagestyle{scrheadings}
\fi
\ifdefined\anspruch@std
  \usepackage{fancyhdr}
  \pagestyle{fancy}
  \fancyhf{}
  \fancyhead[R]{\small\draftmark}
  \fancyhead[L]{\small\nouppercase{\leftmark}}
  \fancyfoot[C]{\thepage}
  \fancypagestyle{plain}{\fancyhf{}\fancyhead[R]{\small\draftmark}%
    \fancyfoot[C]{\thepage}\renewcommand{\headrulewidth}{0pt}}
\fi
\makeatother

% Lists: compact, as in a German Antrag.
\setlist{itemsep=0.2em,topsep=0.4em}

% Tables: a bit more air, ragged-right paragraph columns.
\renewcommand{\arraystretch}{1.15}
\newcolumntype{L}[1]{>{\raggedright\arraybackslash}p{#1}}
 (from docs/dossier/)
% or \input{../preamble} (from docs/dossier/antrag/ or docs/dossier/partner/).
% Each document sets its own \documentclass and its own pdftitle via
% \hypersetup{pdftitle={...}}. Only packages that ship with MiKTeX.

\usepackage[T1]{fontenc}
\usepackage[utf8]{inputenc}
\usepackage{lmodern}
\usepackage[english,ngerman]{babel} % last option = main language: ngerman
\usepackage[a4paper,margin=25mm]{geometry}
\usepackage{booktabs}
\usepackage{longtable}
\usepackage{array}
\usepackage{enumitem}
\usepackage{xcolor}
\usepackage{microtype}
\usepackage[autostyle=true,german=quotes]{csquotes}
\usepackage{marginnote}
\usepackage[hidelinks,unicode,pdflang={de-DE}]{hyperref}

% ---------------------------------------------------------------------------
% Colours
\definecolor{blankbg}{RGB}{255,236,153}   % highlighted placeholder
\definecolor{todofg}{RGB}{170,40,40}      % "to confirm" margin note

% \blank{Text}: a highlighted placeholder "[Text]" for values the sender
% fills in before sending. Usable in running text and in table cells.
\newcommand{\blank}[1]{\begingroup\setlength{\fboxsep}{1pt}%
  \colorbox{blankbg}{\strut[#1]}\endgroup}

% \todo{Text}: a "zu bestätigen" note in the margin. Works in running text,
% lists and longtable cells (via marginnote, not \marginpar).
% Defined with \providecommand so that it does not clash if a document
% loads todonotes on its own (todonotes then wins).
\providecommand{\todo}[1]{\marginnote{\raggedright\scriptsize
  \textcolor{todofg}{\textbf{Zu bestätigen:} #1}}}
\setlength{\marginparwidth}{20mm}

% \draftmark: header text for every page of a draft.
\newcommand{\draftmark}{Entwurf, 23.10.2026}

% Draft header for scrartcl / scrreprt (KOMA) and article / report.
% Letters (scrlttr2, letter) place \draftmark themselves.
\makeatletter
\@ifclassloaded{scrartcl}{\def\anspruch@koma{}}{}
\@ifclassloaded{scrreprt}{\def\anspruch@koma{}}{}
\@ifclassloaded{article}{\def\anspruch@std{}}{}
\@ifclassloaded{report}{\def\anspruch@std{}}{}
\ifdefined\anspruch@koma
  \usepackage[headsepline]{scrlayer-scrpage}
  \clearpairofpagestyles
  \ohead{\small\draftmark}
  \ihead{\small\headmark}
  \cfoot*{\pagemark}
  \automark[section]{section}
  \pagestyle{scrheadings}
\fi
\ifdefined\anspruch@std
  \usepackage{fancyhdr}
  \pagestyle{fancy}
  \fancyhf{}
  \fancyhead[R]{\small\draftmark}
  \fancyhead[L]{\small\nouppercase{\leftmark}}
  \fancyfoot[C]{\thepage}
  \fancypagestyle{plain}{\fancyhf{}\fancyhead[R]{\small\draftmark}%
    \fancyfoot[C]{\thepage}\renewcommand{\headrulewidth}{0pt}}
\fi
\makeatother

% Lists: compact, as in a German Antrag.
\setlist{itemsep=0.2em,topsep=0.4em}

% Tables: a bit more air, ragged-right paragraph columns.
\renewcommand{\arraystretch}{1.15}
\newcolumntype{L}[1]{>{\raggedright\arraybackslash}p{#1}}

\hypersetup{pdftitle={Anspruch: Dossier, Stand 23.10.2026},pdfauthor={Anspruch}}

% Lange URLs dürfen an jeder Stelle umbrechen (nur in diesem Dokument).
\makeatletter
\g@addto@macro\UrlBreaks{\do\a\do\b\do\c\do\d\do\e\do\f\do\g\do\h\do\i\do\j\do\k\do\l
  \do\m\do\n\do\o\do\p\do\q\do\r\do\s\do\t\do\u\do\v\do\w\do\x\do\y\do\z\do\_\do\-
  \do\0\do\1\do\2\do\3\do\4\do\5\do\6\do\7\do\8\do\9}
\makeatother
\Urlmuskip=0mu

% \pfad{...}: Dateipfad in Schreibmaschinenschrift, umbrechbar nach / . und -.
\ExplSyntaxOn
\NewDocumentCommand{\pfad}{m}{
  \tl_set:Nn \l_tmpa_tl {#1}
  \tl_replace_all:Nnn \l_tmpa_tl {/} {/\allowbreak}
  \tl_replace_all:Nnn \l_tmpa_tl {.} {.\allowbreak}
  \tl_replace_all:Nnn \l_tmpa_tl {-} {-\allowbreak}
  \texttt{\l_tmpa_tl}
}
\ExplSyntaxOff

% Spalte wie L aus der Präambel, die auch das erste Wort einer Zelle trennt.
\newcolumntype{P}[1]{>{\raggedright\arraybackslash\hspace{0pt}}p{#1}}

% Kurzzeichen in den Tabellen
\newcommand{\ja}{\textbullet}
\newcommand{\nuetzlich}{$\circ$}
\newcommand{\nein}{--}

\begin{document}

\thispagestyle{plain.scrheadings}
\vspace*{10mm}
{\sffamily\bfseries\Huge Anspruch\par}
\vspace{4mm}
{\sffamily\Large Dossier zum Stand des Projekts, zu Förderung, Partnern und Unterlagen\par}
\vspace{6mm}
{\large Stand: 23.10.2026 \quad\textbullet\quad \draftmark\par}
\vspace{4mm}
Betreiber: \blank{Vor- und Nachname}\\
Kontakt: \blank{E-Mail-Adresse}\\
Adresse der Seite: \blank{Adresse der Seite}
\vspace{6mm}

\begin{minipage}{\textwidth}\small
Dieses Dossier beschreibt, was Anspruch ist, was am 23.10.2026 geprüft und offen ist, wie das
Projekt mit Daten und Recht umgeht, für welche Förderprogramme es in Frage kommt, welche
Organisationen es ansprechen will und welche Unterlagen dafür vorliegen oder fehlen.
Es ist ein Entwurf. Es ist keine Rechtsberatung. Gelb markierte Stellen in eckigen Klammern
sind Lücken, die vor dem Versand ausgefüllt werden. Punkte mit dem Vermerk \enquote{zu
bestätigen} sind nicht an der Quelle geprüft und stehen deshalb nur in den Listen \enquote{Zu
bestätigen}.
\end{minipage}

\vspace{4mm}
\setcounter{tocdepth}{1}
\tableofcontents
\clearpage

% ===========================================================================
\section{Zusammenfassung}

Anspruch ist eine Webseite für Familien, die einen Angehörigen pflegen. Die Familie beantwortet
Fragen zum Alltag der pflegebedürftigen Person. Die Seite schätzt den Pflegegrad, den die
Begutachtung wahrscheinlich ergibt. Sie vergleicht die Leistungen, die daraus folgen, mit dem,
was der Haushalt schon bekommt, und nennt die nächsten Schritte mit ihrer gesetzlichen
Grundlage. Am Ende steht ein PDF auf Deutsch oder Englisch.

\textbf{Das Produkt.} In der Regel genügen 12 bis 20 Fragen statt 64. Steht das Ergebnis noch
nicht fest, zeigt die Seite einen Bereich und nie einen Mittelwert. Jeder Betrag nennt seine
Rechtsquelle und das Prüfdatum. Alles läuft im Browser: kein Konto, kein Server, der Antworten
speichert. Die Oberfläche gibt es in sechs Sprachen, die Fragen selbst auf Deutsch oder
Englisch. Es ist eine Schätzung, keine Begutachtung und keine Rechtsberatung.

\textbf{Stand am 23.10.2026.} Drei Arbeitsläufe am 22.10.2026 haben die Beträge gegen das
geltende Recht geprüft, zwei Zitate berichtigt, eine kritische Sicherheitslücke geschlossen,
vier Fehler aus einem Live-Test behoben, eine abschaltbare Zählung der Nutzung, Partnerlinks,
Impressum und Datenschutzerklärung gebaut. Offen sind vor allem: die Prüfung der vier
Fremdsprachen durch Muttersprachler, das Pflegeneuordnungsgesetz (PNOG, Kabinettsentwurf vom
30.09.2026), drei Fragen zum Quelltext und eine Quelle mit Prüfdatum 26.08.2026.

\textbf{Datenschutz.} Was die Familie tippt, verlässt das Gerät nicht. Die Zählung ist
standardmäßig aus. Ist sie an, zählt Plausible (nach eigenen Angaben in der EU) Seitenaufrufe
und fünf Schritte ohne Cookies und ohne Antworten. Vor dem Einschalten stehen neun
Entscheidungen aus, darunter zwei Auftragsverarbeitungsverträge und eine kurze rechtliche
Prüfung.

\textbf{Förderung.} Der Prototype Fund (Frist 30.11.2026) ist das einzige Programm ausdrücklich
für eine Privatperson; die Passung ist vorher zu erfragen. Modellvorhaben nach § 123 SGB XI
laufen über ein Land oder eine Kommune, meist mit einem Träger; nur Niedersachsen lässt auch
natürliche Personen zu. Aktion Mensch verlangt einen gemeinnützigen Träger, die Deutsche
Fernsehlotterie eine Trägerorganisation. Die Robert Bosch Stiftung nimmt keine
unaufgeforderten Anträge an.

\textbf{Partner.} Für einen Pilotversuch empfehlen wir als erste drei: einen Pflegestützpunkt in
\blank{Land}, eine Migrationsberatung (MBE) eines Wohlfahrtsverbands und eine Beratungsstelle
des Sozialverbands VdK.

\textbf{Unterlagen.} Impressum, Datenschutzerklärung, Compliance-Liste, Pilotvorschlag und
Förderübersicht liegen im Repository. Sieben Entwürfe entstehen in diesem Lauf. Es fehlen
Lizenzdatei, Verzeichnis der Verarbeitungstätigkeiten, Datenschutzkonzept,
Barrierefreiheitserklärung und Lebenslauf; die Auftragsverarbeitungsverträge sind
abzuschließen.

\textbf{Die nächsten 90 Tage.} Lizenz entscheiden und Fragen an den Prototype Fund stellen,
Rechtsseiten mit ausgefüllten Lücken live schalten, Verträge abschließen, bis 30.11.2026 die
Bewerbung einreichen, drei Partner ansprechen, einen Pilot starten.

\subsection*{Vor dem Versand ausfüllen}
{\small
\begin{itemize}[nosep]
  \item \blank{Vor- und Nachname}, \blank{E-Mail-Adresse}: Betreiber und Kontakt; auch in
        \pfad{lib/legal/blanks.ts}.
  \item \blank{Adresse der Seite}: Produktionsadresse; auch in \pfad{docs/DEPLOY.md}.
  \item \blank{Land}: Bundesland des Wohnsitzes; bestimmt Aufsichtsbehörde und Weg nach § 123.
  \item \blank{Kommune oder Träger}: Antragsteller für § 123, Aktion Mensch, Fernsehlotterie.
  \item \blank{Team}: Namen und Rollen für den Prototype Fund.
  \item \blank{Betrag}: Preis eines Pilotversuchs oder \enquote{kostenlos}.
  \item \blank{Startdatum}: Beginn des ersten Pilotversuchs.
\end{itemize}
}

\clearpage
% ---------------------------------------------------------------------------
\begin{otherlanguage}{english}
\addsec{Summary}

Anspruch is a website for families who care for a relative. The family answers questions about
the everyday life of the person who needs care. The site estimates the Pflegegrad the
assessment would likely arrive at. It compares the benefits that follow with what the household
already receives, and lists the next steps with their statutory basis. At the end it produces a
PDF in German or English.

\textbf{The product.} Typically 12 to 20 questions are enough instead of 64. Where the result is
not settled, the site shows a range and never a midpoint. Every amount names its legal source
and the date it was checked. Everything runs in the browser: no account, no server that stores
answers. The interface is in six languages, the questions themselves in German or English. It
is an estimate, not an assessment and not legal advice.

\textbf{State on 23 October 2026.} Three work runs on 22 October 2026 checked the amounts
against current law, corrected two citations, closed a critical security advisory, fixed four
faults from a live test, and built switchable usage counting, partner links, an Impressum and a
privacy policy. Still open, above all: a native-speaker review of the four other languages, the
Pflegeneuordnungsgesetz (PNOG, cabinet bill of 30 September 2026), three questions about the
source text, and one source last checked on 26 August 2026.

\textbf{Privacy.} What the family types does not leave the device. Counting is off by default.
When on, Plausible (in the EU, by its own account) counts page views and five steps without
cookies and without answers. Nine decisions are pending before it is switched on, among them
two data processing agreements and a short legal read.

\textbf{Funding.} The Prototype Fund (deadline 30 November 2026) is the only programme
explicitly for a private individual; its fit must be asked first. Model projects under § 123
SGB XI run through a Land or a Kommune, usually with a Träger; only Niedersachsen also admits
natural persons. Aktion Mensch requires a non-profit Träger, the Deutsche Fernsehlotterie a
Trägerorganisation. The Robert Bosch Stiftung does not accept unsolicited applications.

\textbf{Partners.} For a pilot we recommend these first three: a Pflegestützpunkt in
\blank{Land}, a migration counselling service (MBE) of a welfare association, and an advice
centre of the Sozialverband VdK.

\textbf{Documents.} Impressum, privacy policy, compliance checklist, pilot proposal and funding
outline are in the repository. Seven drafts are written in this run. Missing: a licence file,
a record of processing activities, a data protection concept, an accessibility statement and a
CV; the data processing agreements are still to be concluded.

\textbf{The next 90 days.} Decide the licence and put the questions to the Prototype Fund, put
the legal pages live with their blanks filled, conclude the agreements, submit the application
by 30 November 2026, approach three partners, start a pilot.

The blanks to fill before sending are listed at the end of the German summary.
\end{otherlanguage}

\clearpage
% ===========================================================================
\section{Das Produkt}

Quelle dieses Abschnitts ist \pfad{README.md}.

\subsection{Was es tut}
Die Familie beantwortet Fragen zum Alltag. Die Seite schätzt den Pflegegrad, den die
Begutachtung nach dem Neuen Begutachtungsassessment wahrscheinlich ergibt. Sie vergleicht den
daraus folgenden Anspruch mit dem, was der Haushalt schon bekommt. Sie nennt, was zu tun ist,
in Reihenfolge und mit der Rechtsgrundlage jedes Schritts. Am Ende steht ein PDF.

Die Seite sagt auf jedem Weg: Dies ist eine Schätzung, keine Begutachtung durch den
Medizinischen Dienst. Sie ist keine Rechtsberatung. Die Pflegeberatung nach § 7a SGB XI ist
kostenlos und ein Rechtsanspruch.

\subsection{Vier Grundsätze}
\begin{description}
  \item[Kurze Befragung.] In der Regel 12 bis 20 Fragen statt 64. Die Befragung endet, wenn
    weitere Antworten das Ergebnis nachweislich nicht mehr ändern können. Dafür wird für jede
    offene Frage der günstigste und der ungünstigste Fall gerechnet. Zwei Vorfragen schließen
    die Module 2 und 3, Modul 5 wird auf der Stufe gefragt, auf der es bewertet wird, und
    Kriterien einer Tätigkeit werden zusammen gefragt. Nur das Zusammenfassen kann Genauigkeit
    kosten; jede Gruppe lässt sich wieder aufklappen. Der Bericht sagt, wie viele der 64 Fragen
    gestellt wurden und warum die übrigen entfallen sind.
  \item[Bereich statt Mittelwert.] Stimmen die beiden Grenzen nicht überein, zeigt der Bericht
    einen Bereich, zum Beispiel \enquote{zwischen Pflegegrad 2 und 4}, und bietet die offenen
    Fragen an. Er wählt nie die Mitte.
  \item[Nachvollziehbare Beträge.] Jeder Betrag und jede Schwelle trägt eine Quelle mit dem
    Datum der letzten Prüfung (\pfad{lib/rules/sources.ts}). Ein Betrag ohne Quelle wird
    nicht gezeigt. Einmalige Zuschüsse zählen nicht zur Monatssumme. Bedingte Leistungen
    erscheinen als \enquote{prüfen, ob das zutrifft}. Pflegegeld und Pflegesachleistung werden
    nicht addiert.
  \item[Auf dem Gerät.] Alles läuft im Browser. Es gibt kein Konto und keinen Server, der
    Antworten speichert. Die Antworten bleiben in diesem Browser, damit eine halb fertige
    Befragung einen geschlossenen Tab übersteht; beim nächsten Besuch kann man fortfahren oder
    alles löschen. Geschrieben wird erst, wenn eine Frage beantwortet ist.
\end{description}

\subsection{Sprachen}
Die Oberfläche gibt es auf Deutsch, Englisch, Türkisch, Russisch, Polnisch und Arabisch.
Arabisch läuft von rechts nach links. Die Fragen selbst und der Bericht bleiben auf Deutsch oder
Englisch, weil sie ein Rechtsinstrument wiedergeben und eine ungeprüfte Übersetzung geglaubt
würde. Die Oberfläche sagt das in der Sprache der Leserin und lässt sie wählen, in welcher der
beiden Sprachen die Fragen erscheinen. Die Fragen sind in Alltagssprache formuliert; der amtliche
Wortlaut steht klein darunter.

\subsection{Barrierefreiheit}
Schriftgröße, Farbschema, Sprache und Tempo werden auf dem ersten Bildschirm eingestellt. Es
gibt vier Farbschemata; jedes Textpaar erreicht WCAG AAA (7:1), durch Tests festgeschrieben.
Fragen können vorgelesen und durch Nennen der Nummer beantwortet werden; Antippen geht immer.
Spracheingabe braucht einen Browser, der sie kann (Chrome, Edge). Standard ist eine Frage pro
Bildschirm mit Schaltflächen von mindestens 44 Pixeln.

\subsection{Das PDF}
Das PDF wird im Browser erzeugt und nirgendwohin gesendet. Es ist auf Deutsch oder Englisch,
weil es bei der Pflegekasse über den Tisch geht. Es enthält die Schätzung, die Beträge, die
Schritte mit Paragraf und die gegebenen Antworten, damit die Familie bei der Begutachtung nicht
aus dem Gedächtnis antworten muss.

\subsection{Bekannte Grenzen}
Die Fremdsprachen sind nicht von Muttersprachlern geprüft. Vorlesen einer Frage dauert etwa
20~Sekunden. Die maschinelle Übersetzung des Browsers ist abgeschaltet. Bei zusammengefassten
Fragen muss eine Familie die Gruppe selbst aufklappen, wenn sich die Kriterien unterscheiden.
Der Quellcode hat noch keine Lizenz.

% ===========================================================================
\section{Stand am 23.10.2026}

\subsection{Geprüft und behoben}
Quellen: die Zusammenfassungen der drei Läufe vom 22.10.2026 in
\pfad{.orchestration/runs/} (LC = \pfad{2026-10-22-last-check}, LT =
\pfad{2026-10-22-live-test-fixes}, UP = \pfad{2026-10-22-usage-and-partners}; Zahlen nach
dem Doppelpunkt sind Zeilen).

{\small\setlength{\tabcolsep}{3pt}
\begin{longtable}{@{}P{17mm}P{95mm}P{40mm}@{}}
\toprule
\textbf{Datum} & \textbf{Was} & \textbf{Beleg}\\
\midrule
\endhead
22.10.2026 & Oberfläche auf Türkisch, Russisch, Polnisch und Arabisch Schlüssel für Schlüssel gegen
  das Deutsche geprüft (durch Sprachmodelle); sinnverändernde Fehler behoben, darunter fehlende
  Fachbegriffe \enquote{Widerspruch} und \enquote{Pflegeberatung}. Eine Überschrift in allen
  sechs Sprachen berichtigt (fünf statt vier Vorfragen). & LC:13, LC:22; \pfad{lib/i18n/strings.ts} \\
22.10.2026 & Alle Beträge, Schwellen und Fristen nach geltendem Recht richtig. Zwei Zitate
  berichtigt: Wohngruppenzuschlag nach § 45f statt § 38a SGB XI; die Zwei-Wochen-Regel nach
  § 7b Abs.~1 statt § 7a. Prüfdaten auf 22.10.2026 erneuert, außer einem (siehe Offene Punkte, 3). & LC:14, LC:24;
  \pfad{lib/rules/sources.ts} \\
22.10.2026 & Kritische Sicherheitslücke in Next.js geschlossen (Wechsel auf 16.4.2);
  \texttt{npm run typecheck} läuft in einem frischen Klon. & LC:15--16, LC:28 \\
22.10.2026 & Prüfung: Typen, 293 Tests, Lint, Produktions-Build, Durchlauf im Browser auf
  Deutsch und Arabisch, PDF auf dem Gerät. In \texttt{main} zusammengeführt (PR \#6). & LC:18;
  \texttt{git log main} \\
22.10.2026 & Haushalt mit unbrauchbaren Armen und Beinen: Befragung läuft weiter, bis der
  Punktegrad feststeht; Ergebnis Pflegegrad 5 mit dem Punktegrad daneben. & LT:9 \\
22.10.2026 & Hinweis bei Modul 5 neu betitelt; Fristen als Datum im Format der Sprache;
  Situationsbildschirm \enquote{Vorab: Ihre Situation} statt \enquote{Schritt 0 von 19}. & LT:10 \\
22.10.2026 & Schaltfläche \enquote{N weitere Fragen beantworten} wirkt jetzt auch in der
  Listenansicht. Prüfung: 313 Tests, Typen, Lint, Build. & LT:11, LT:13 \\
22.10.2026 & Zählung der Nutzung, standardmäßig aus; Partnerlinks mit \texttt{?via=};
  \pfad{docs/DEPLOY.md}, \pfad{docs/PARTNERS.md}. & UP:13 \\
22.10.2026 & Pilotvorschlag und Förderübersicht (\pfad{docs/proposal-organisations.md},
  \pfad{docs/funding-outline.md}). & UP:14 \\
22.10.2026 & Seiten \pfad{/impressum} und \pfad{/datenschutz} (Deutsch und Englisch),
  Fußzeilen-Links, \pfad{docs/COMPLIANCE.md}. Prüfung: 346 Tests, Typen, Lint, Build; keine
  Anfrage erreichte einen Anbieter. & UP:15, UP:17 \\
\bottomrule
\end{longtable}
}

Die Änderungen aus den Läufen LT und UP liegen auf den Branches \texttt{live-test-fixes}
(PR \#7) und \texttt{usage-and-partners} (PR \#8); \texttt{main} enthält sie noch nicht. Die
Seite ist auf Vercel bereitgestellt (UP:7); ob die bereitgestellte Fassung Impressum und
Datenschutzerklärung schon enthält, ist zu prüfen.

\subsection{Offene Punkte}
\begin{enumerate}
  \item \textbf{Prüfung durch Muttersprachler.} Je eine Person für Türkisch, Russisch,
    Polnisch und Arabisch, die das deutsche Pflegesystem kennt, soll zuerst die fünf Texte mit
    rechtlichem Inhalt lesen (\texttt{privacyBody}, \texttt{bescheidHint}, \texttt{m5CapBody},
    \texttt{exclusiveNote}, \texttt{disclaimerExtra}), dazu die neuen Texte
    \texttt{m5SkippedTitle}, \texttt{situationStep}, \texttt{pointsRangeTitle},
    \texttt{gradeNoneNominative} und die Texte der Zählung. Offen sind auch Stilfragen: Polnisch
    duzt, wo das Deutsche siezt; arabische Schaltflächen mischen Formen; ob Pflegegeld,
    Pflegesachleistung und Pflegekasse beim ersten Auftreten erklärt werden. (LC:32--33, LT:21,
    UP:31)
  \item \textbf{PNOG.} Der Kabinettsentwurf des Pflegeneuordnungsgesetzes vom 30.09.2026 würde
    die meisten Beträge ab 01.01.2027 ersetzen und die Schwellen für Pflegegrad 1 bis 3 auf
    15/30/50 Punkte senken. Die erste Lesung ist für den 05.11.2026 angesetzt. Nichts davon
    erreicht die Nutzer; die Hinweise stehen in \pfad{lib/rules/sources.ts}. Wird es wie
    entworfen beschlossen, sind die Beträge der Seite ab 01.01.2027 falsch. (LC:26, LC:36)
  \item \textbf{Bedarfskonstellation.} Die Seite zeigt als Stand des Regelwerks den 26.08.2026,
    weil die Angabe, die Begutachtungs-Richtlinien kennten genau eine besondere
    Bedarfskonstellation, nicht neu abgerufen wurde. (LC:35)
  \item \textbf{Drei Fragen zum Quelltext}, auch für Deutsch und Englisch:
    \texttt{pdfLanguageNote} sagt, die Pflegekasse lese das PDF auf Deutsch oder Englisch; ein
    englisches PDF ist fraglich. \texttt{exclusiveNote} sagt, Pflegegeld und
    Pflegesachleistung seien \enquote{keine Summe}, § 38 SGB XI erlaubt aber eine anteilige
    Kombination. \texttt{estimateHigherBody} nennt \enquote{29 von 100} ohne Quelle, gegen die
    eigene Regel des Projekts; dieselbe Zahl steht in \pfad{README.md}. (LC:34)
  \item \textbf{Lizenz.} \enquote{Not yet chosen} (\pfad{README.md}). Empfehlung: EUPL-1.2 mit
    alleinigem Urheberrecht (\pfad{docs/funding-outline.md}).
  \item \textbf{Kleinere technische Punkte.} Ein Bereich mit unterer Grenze null erscheint als
    \enquote{Zwischen Pflegegrad 0 und Pflegegrad 3}; \pfad{lib/rules/gap.ts:325} vergleicht
    die Widerspruchsfrist als Mitternacht UTC; eine verbleibende Warnung zu \texttt{braces} nur
    zur Lint-Zeit. (LT:22--24, LC:28)
\end{enumerate}

% ===========================================================================
\section{Datenschutz und Compliance}

Quelle: \pfad{docs/COMPLIANCE.md}, in dessen Reihenfolge; dort mit Zeilenangaben der
Rechercheberichte (SC, SA). Entwurf für eine rechtliche Prüfung, keine Rechtsberatung.

\subsection{Haltung}
\begin{itemize}
  \item \textbf{Nichts verlässt das Gerät.} Antworten bleiben in diesem Browser und lassen sich
    mit einem Tipp löschen. Kein Konto, kein Server, der Antworten speichert. Die App setzt
    keine Cookies. Ausnahme: Bei Spracheingabe erkennt der Browser die Sprache, möglicherweise
    auf Servern des Browserherstellers; der Betreiber erhält weder Ton noch Text.
  \item \textbf{Zählung standardmäßig aus.} Ohne die Umgebungsvariablen enthält der Build kein
    Skript und sendet nichts. Ist sie an, zählt Plausible Seitenaufrufe (ohne Query-String) und
    fünf Ereignisse: \texttt{intake\_started}, \texttt{result\_shown},
    \texttt{questions\_refined}, \texttt{pdf\_downloaded}, \texttt{intake\_erased}. Jedes trägt
    Oberflächensprache, Tempo und Partnerkürzel; \texttt{result\_shown} dazu Spannen für die
    Zahl der Fragen, die Art des Ergebnisses (keins, Bereich, Grad) und die monatliche Lücke;
    \texttt{pdf\_downloaded} die Sprache des PDF. Nie gesendet: Grad, Punkte, Antworten, Datum,
    Freitext, Kennung, Cookie. Eine Erlaubnisliste im Code und ein Test sichern das.
    (\pfad{README.md})
  \item \textbf{Partnerlink nur im Speicher.} Das Kürzel aus \texttt{?via=} wird geprüft, sofort
    aus der Adresszeile entfernt und nur für diesen Seitenbesuch im Arbeitsspeicher gehalten;
    nie auf dem Gerät gespeichert. (\pfad{docs/PARTNERS.md})
  \item \textbf{Rechtsseiten.} \pfad{/impressum} (§ 5 DDG) und \pfad{/datenschutz}, Deutsch
    und Englisch, Texte in \pfad{lib/legal/}, mit markierten Lücken.
  \item \textbf{Hosting bei Vercel.} Vercel Inc., USA, verarbeitet als Auftragsverarbeiter
    IP-Adresse, ungefähren Ort und Serverprotokolle. Vercel sagt keine reine EU-Verarbeitung
    zu; die Datenschutzerklärung nennt das EU-US Data Privacy Framework und
    Standardvertragsklauseln. (\pfad{lib/legal/datenschutz.ts}, UP:22)
  \item \textbf{Plausible in der EU.} Plausible Insights OÜ, Estland; nach eigenen Angaben
    Verarbeitung und Speicherung in der EU, bei Hetzner in Deutschland. Das ist eine Angabe des
    Anbieters. (\pfad{lib/legal/datenschutz.ts})
\end{itemize}

\subsection{Lücken (COMPLIANCE, Abschnitt 1)}
Alle Lücken stehen an einer Stelle, \pfad{lib/legal/blanks.ts}; ein Test schlägt fehl, wenn
eine hinzukommt oder fehlt. Zu füllen: Vor- und Nachname; Straße, Hausnummer, Postleitzahl und
Ort (kein Postfach); E-Mail-Adresse; Umsatzsteuer- oder Wirtschafts-ID, Register,
Aufsichtsbehörde, Berufsangaben nur, wenn sie zutreffen; die Datenschutz-Aufsichtsbehörde des
Landes.

\subsection{Ausstehende Entscheidungen (COMPLIANCE, Abschnitt 2)}
{\small\setlength{\tabcolsep}{3pt}
\begin{longtable}{@{}P{6mm}P{38mm}P{89mm}P{17mm}@{}}
\toprule
\textbf{Nr.} & \textbf{Entscheidung} & \textbf{Grundlage} & \textbf{Beleg}\\
\midrule
\endhead
1 & Land und Aufsichtsbehörde & Für eine Privatperson wohl die Behörde des Wohnsitzlandes (zu
  bestätigen). & SC:52 \\
2 & Vercel Hobby oder Pro & Hobby nur für nichtkommerzielle private Nutzung; Partnerlinks und
  Zählung für Organisationen können als kommerziell gelten (zu bestätigen). Der Plan bestimmt
  die Aufbewahrung der Laufzeitprotokolle (Hobby 1 Stunde, Pro 1 Tag); die Seite nennt sie
  noch nicht. & SC:18, SC:48 \\
3 & Plausible-Tarif & Eigene Eigenschaften (Partnerkürzel, Spannen) brauchen den
  Business-Tarif; Preise und Grenzen bei Anmeldung prüfen. & SA:6--7, 43, 51 \\
4 & AVV mit Vercel & Der Vertrag ist \enquote{Teil des Enterprise Agreement}, ToS §~11.1 bezieht
  ihn angeblich für alle ein; bestätigen, dass er den gewählten Plan abdeckt. Unterauftragnehmer
  zu prüfen. & SC:20, 49, 59 \\
5 & AVV mit Plausible & Durch Nutzung angenommen, ohne Unterschrift; nur EU, keine
  Standardvertragsklauseln. Unterauftragnehmer außer Hetzner und Bunny zu prüfen. & SC:31 \\
6 & Rechtliche Prüfung DSK OH Digitale Dienste Rn.~87--90 & Keine pauschale Einwilligungsfreiheit
  für Reichweitenmessung; Ereignisse mit Eigenschaften gehen über eine reine Zählung hinaus.
  Die Seite stützt sich auf Art.~6 Abs.~1 lit.~f DSGVO und sagt nicht, dass keine Einwilligung
  nötig sei. Auch Rn.~68/69 prüfen (zu bestätigen): eine Abwägung nach lit.~f allein erfüllt
  § 25 Abs.~2 Nr.~2 möglicherweise nicht. & SC:26, 60; SA:32--36, 40 \\
7 & Rechtliche Prüfung der Speicherung im Browser (Rn.~80) & Das Aufbewahren der Antworten über
  Besuche hinweg ist der weiche Punkt; außerhalb von Login-Diensten ist Speicherung nur für die
  Sitzung der Regelfall. Vorhanden: nichts wird vor einer Antwort geschrieben, Löschen mit einem
  Tipp. & SC:24--25 \\
8 & Verzeichnis nach Art.~30 & Ein einseitiges VVT mit drei Einträgen: Vercel-Protokolle,
  Plausible-Statistik, E-Mail-Kontakt aus dem Impressum. Die Ausnahme nach Art.~30 Abs.~5
  greift wahrscheinlich nicht (zu bestätigen). & SC:43 \\
9 & BFSG und BITV & Ein kostenloser Schätzer ohne Vertrag fällt wohl nicht unter das BFSG (zu
  bestätigen, § 2 Nr.~26). Öffentliche Stellen, die ihn einbinden, fallen unter BITV 2.0;
  Barrierefreiheitserklärung auf Anfrage anbieten. & SC:55--56 \\
\bottomrule
\end{longtable}
}

\subsection{Vor dem Einschalten der Zählung (COMPLIANCE, Abschnitt 3)}
In dieser Reihenfolge: (1) alle Lücken füllen, Tests grün; (2) Land und Behörde in
\texttt{DPA\_AUTHORITY}; (3) Vercel-Plan wählen, Protokollaufbewahrung gegebenenfalls in
Abschnitt 5 der Datenschutzerklärung; (4) beide AVV bestätigen und ablegen; (5) auf der
Produktionsadresse mit \texttt{curl -I} auf Cookies prüfen; (6) Aufbewahrungsdauer bei
Plausible bestätigen und in Abschnitt 6 eintragen; (7) entscheiden, wie ein Widerspruch nach
Art.~21 umgesetzt wird; (8) bestätigen, dass \texttt{?via=} Plausible nicht in einer URL
erreicht; (9) EU-Datenregion im Plausible-Dashboard bestätigen; (10) rechtliche Prüfung (6, 7)
einholen und VVT schreiben; (11) \pfad{/impressum} und \pfad{/datenschutz} mit gefüllten
Lücken live; (12) erst dann \texttt{NEXT\_PUBLIC\_ANALYTICS\_PROVIDER} setzen.

\subsection{Offene Punkte aus dem Code (COMPLIANCE, Abschnitt 4)}
\begin{itemize}
  \item \textbf{Spracheingabe.} Nutzt die Spracherkennung des Browsers
    (\pfad{lib/a11y/speech.ts}). Die Datenschutzerklärung sagt: freiwillig, vom Browser
    ausgeführt, möglicherweise auf Servern des Browserherstellers, der Betreiber erhält weder Ton
    noch Text; sie nennt keinen Anbieter. Zu tun: prüfen, wie Chrome und Edge Sprache
    verarbeiten, und ob sie genannt werden.
  \item \textbf{Spanne der Lücke.} \texttt{result\_shown} sendet die monatliche Lücke in fünf
    Spannen. Sie ist aus dem Ergebnis abgeleitet, keine Antwort und kein Grad; in der
    rechtlichen Prüfung (6) bestätigen lassen, dass das ohne Einwilligung zulässig ist.
\end{itemize}

\subsection{Was bewusst nicht geschieht (COMPLIANCE, Abschnitt 5)}
Keine Cookies der App, nichts für die Zählung auf dem Gerät. Keine Konten; die App sendet
keine Antwort (Spracheingabe siehe oben). Kein Grad, keine Punkte, keine Antwort in der
Zählung. Kein Partnerkürzel im Speicher des Geräts. Keine externen Schriften. Keine Werbung,
keine Tracker Dritter außer Plausible, wenn eingeschaltet. Kein Cookie-Banner, vorbehaltlich
der rechtlichen Prüfung (6).

% ===========================================================================
\section{Förderprogramme}

Quellen: \pfad{scout-applications.md} (SA, Rechercheprüfung vom 23.10.2026, mit
Zeilennummer und Kennzeichen V = an der Quelle geprüft, A = archivierte Fassung, S = Sekundär,
TV = zu prüfen) und \pfad{docs/funding-outline.md} (FO). In der Tabelle steht nur, was V ist
oder aus FO stammt; alles andere steht unter \enquote{zu bestätigen}.

{\small\setlength{\tabcolsep}{3pt}
\begin{longtable}{@{}P{18mm}P{19mm}P{28mm}P{25mm}P{23mm}P{17mm}P{14mm}@{}}
\toprule
\textbf{Programm} & \textbf{Geldgeber} & \textbf{Passung} & \textbf{Berechtigt} &
\textbf{Betrag} & \textbf{Frist} & \textbf{Status}\\
\midrule
\endhead
Prototype Fund, Jahrgang 03 & zu bestätigen & Unsicher; Linien Resilienz (\enquote{Freie
  Alternativen zu proprietärer, zentralisierter \dots\ Software}) oder Innovation
  (\enquote{Datensicherheit und Software-Infrastruktur}) & zu bestätigen & zu bestätigen &
  30.11.2026 & erst fragen \\
§ 123 SGB XI, Modellvorhaben & Pflegeversicherung (Ausgleichsfonds) und Land oder Kommune
  (hälftig) & Anspruch als ein Baustein; \enquote{Digitale Vernetzung der Pflegeangebote} ist
  vorgesehen, ein Schätzer nicht genannt & je Land; Antrag an Land oder Kommune & bis
  30 Mio. EUR im Jahr aus der Pflegeversicherung, mit Kofinanzierung bis 60 Mio. im Jahr, 2025 bis 2028 & je Land &
  braucht Träger \\
§ 123 in Niedersachsen & Land (45\,\% Kommunen, 50\,\% andere) und Pflegeversicherung &
  wie oben, Durchführung in Niedersachsen & natürliche und juristische Personen & bis 1~Mio.~EUR
  je Projekt & jährlich 31.08. & erst fragen \\
§ 123 in Hessen & Land und Pflegeversicherung & Ergebnis darf nicht kommerziell vertrieben werden
  & schriftlicher Antrag mit Stellungnahme des Kreises & mindestens 30.000 EUR im Jahr; bis
  31.12.2028 & zu bestätigen & braucht Träger \\
§ 123 in Berlin & Land und Pflegeversicherung & wie oben & nur juristische Personen &
  200.000 EUR im Jahr für alle Projekte & 21.08.2026, abgelaufen & braucht Träger \\
Aktion Mensch & Aktion Mensch & Kein Programm zu digitaler Teilhabe; neue Programme ab
  02.11.2026 & gemeinnützig, Sitz in Deutschland; keine Einzelpersonen & zu bestätigen &
  Antragsphase ab 01.01.2027 & braucht Träger \\
Deutsche Fernsehlotterie, Stiftung Deutsches Hilfswerk & Stiftung Deutsches Hilfswerk &
  Passung zu bestätigen & eine Bewerbung je Trägerorganisation; nur ohne Hilfswerk-Förderung
  2026 & Gesamtkosten mindestens 50.000~EUR & Daten der Organisation bis 10.12.2026; Bewerbung
  02.11.--17.12.2026 & braucht Träger \\
Robert Bosch Stiftung & Robert Bosch Stiftung & Keine; nur \enquote{Deutscher Schulpreis
  2027} offen & keine unaufgeforderten Anfragen & \nein & \nein & entfällt \\
Google.org & \nein & \nein & \nein & \nein & \nein & entfällt \\
\bottomrule
\end{longtable}
}

Belege je Zeile: Prototype Fund SA:7 (Frist, Linien), FO Tabelle und Frage 4 (Passung
unsicher, erst fragen). § 123 SA:19 und FO Tabelle. Niedersachsen SA:21. Hessen SA:22, SA:50.
Berlin SA:20. Aktion Mensch SA:27--29. Fernsehlotterie SA:34. Robert Bosch Stiftung SA:35.
Google.org FO Frage 16 (\enquote{Drop unless a call for Germany appears in 2027}).

\subsection{Prototype Fund, Jahrgang 03}
\textbf{Geprüft (SA:7).} Drei Linien auf \url{https://bewerben.prototypefund.de/}, jeweils mit
Frist 30.11.2026: Up and Coming (unter 26, im Studium oder in einer informatischen Ausbildung), Resilienz, Innovation.
Die Seite nennt keine Beträge, keine Berechtigung, keine Lizenzliste und keinen Geldgeber. Das
Formular ist nur nach Anmeldung lesbar.

\textbf{Nach dem Handbuch Jahrgang 02, archiviert (A), gegen J03 zu prüfen.} Fragen und
Grenzen (SA:10): Allgemeine Informationen; Projektbeschreibung 100 Wörter; Gesellschaftliche
Herausforderung 175; Technische Umsetzung 175; Aktueller Stand 100; Link (optional); Innovation
und ähnliche Ansätze 100; Zielgruppe 100; Meilensteine 100 (\enquote{1/ Backend, 2/ Frontend}
gilt als zu grob); Team 30; Erfahrung, höchstens drei Projekte, 100; Arbeitsstunden (Zahl);
Motivation 100; Second Stage ja oder nein, Schwerpunkt, Meilensteine; Erklärungen und
Checkliste. Keine Anhänge, nur Text und Links. Ein öffentliches Repository ist nicht Pflicht,
hilft aber (SA:15). Zustimmung zur Veröffentlichung unter einer Open-Source-Lizenz (SA:14).

\textbf{Unterlagen.} Projektbeschreibung und Zeitplan aus diesem Lauf liefern die Texte für die
Formularfelder (\pfad{antrag/projektbeschreibung.tex}, \pfad{antrag/zeitplan.tex}); dazu
die Lizenzentscheidung und die Angaben zum Team (\blank{Team}).

\textbf{Vorgehen.} Erst bewerben, wenn die Fragen zu Berechtigung, Passung und Lizenz beantwortet
sind (FO). Wird ein Träger Antragsteller, ist dieser Weg verschlossen (SA:40).

\subsection{Modellvorhaben nach § 123 SGB XI}
\textbf{Geprüft (SA:19).} Grundlage sind die Empfehlungen des GKV-Spitzenverbandes nach
§ 123 Abs.~3 SGB XI (beschlossen 04.11.2024, aktualisiert 20.05.2026). Hälftige
Kofinanzierung durch Land und/oder Kommune; höchstens vier Jahre, Ende spätestens 31.12.2029;
keine Anschlussförderung. Voraussetzungen: Antrag an Land oder Kommune vor Beginn; eine
Konzeption; wissenschaftliche Begleitung und Auswertung nach § 124; Förderzusage der
Landesverbände der Pflegekassen und des PKV-Verbands; Einvernehmen des BMG.

\textbf{Antrag in Textform (Nr.~5.1):} Name und Anschrift des Modellträgers; Beschreibung des
Inhalts; örtlicher Geltungsbereich; ob bei anderen Geldgebern Anträge gestellt sind; Angaben zur
wissenschaftlichen Begleitung. \textbf{Konzeption (Nr.~2.2):} Ziele, Inhalt, Beginn und Dauer,
Umsetzung, Kosten für Projekt und Auswertung (Finanzierungsplan), regional innovativer Charakter.

\textbf{Länder (V).} Niedersachsen lässt als einziges gesehenes Land natürliche Personen zu;
Formulare: Richtlinie, Antragsvordruck, Personalbogen, elektronisch; zuständig ist das
Landesamt für Soziales, Jugend und Familie, Außenstelle Osnabrück (SA:21). Hessen verlangt eine
Stellungnahme des Landkreises oder der kreisfreien Stadt (Anlage~3), die Auswertung nach § 124,
Personalkosten höchstens nach TV-H/TVöD, und das Ergebnis darf nicht kommerziell vertrieben
werden (SA:22). Damit wäre der Weg einer zusätzlichen kostenpflichtigen Lizenz (FO) für ein in
Hessen gefördertes Ergebnis verschlossen (SA:50). Berlin: nur juristische Personen,
zweistufig, die Frist 21.08.2026 ist abgelaufen (SA:20). Schleswig-Holstein: Richtlinie vom
01.10.2025, Land 1~Mio.~EUR im Jahr, Antragsteller nicht genannt (SA:23).

\textbf{Unterlagen.} Konzeption: \pfad{antrag/projektbeschreibung.tex}; Finanzierungsplan:
\pfad{antrag/finanzplan.tex}; Zeitplan: \pfad{antrag/zeitplan.tex}; Partnerzusagen:
\pfad{partner/letter-of-intent.tex}. Antragsteller: \blank{Kommune oder Träger} in
\blank{Land}.

\subsection{Aktion Mensch und Deutsche Fernsehlotterie, über einen Träger}
\textbf{Aktion Mensch (V, SA:27--29).} Gemeinnützige Organisationen mit Sitz in Deutschland
(e.V., gGmbH, gUG, gAG, KdöR, eG, rechtsfähige Stiftung); \enquote{Einzelpersonen sind leider
nicht förderfähig}; öffentlicher oder kommerzieller Einfluss unter 50\,\%; Nachweis durch einen
Bescheid des Finanzamts. Neue Programme ab 02.11.2026, nächste Antragsphase ab 01.01.2027.
Vorlage: der Bescheid des Finanzamts; die übrigen Unterlagen stammen aus älteren Quellen und
stehen unter \enquote{zu bestätigen}.

\textbf{Deutsche Fernsehlotterie (V, SA:34).} Organisationsdaten bis 10.12.2026,
Fördermittelbewerbung 02.11.--17.12.2026, Entscheidung Mai 2027, frühester Projektbeginn
01.06.2027. Eine Bewerbung je Trägerorganisation; nur Organisationen ohne Förderung des
Hilfswerks in den Runden 2026; Gesamtkosten mindestens 50.000~EUR.

\subsection{Die übliche Antragsmappe}
Der \enquote{Kleine Leitfaden für Förderanfragen an Stiftungen} (Bundesverband Deutscher
Stiftungen; V, SA:43) nennt: Anschreiben mit beantragtem Betrag; Titelseite; Zusammenfassung
(bis eine Seite); Langfassung mit Einleitung und Ziel, Bedarf, Zielerreichung, Zeit- und
Maßnahmenplan, Personal, Dokumentation und Auswertung, Kosten- und Finanzierungsplan,
Perspektive; Anhang mit Jahresbericht, Organigramm, Referenzen oder Empfehlungsschreiben.
Unterschrift durch Geschäftsführung oder Vorstand. Freistellungsbescheid, Registerauszug und
Satzung verlangen die Formulare einzelner Geldgeber. Für Datenschutzkonzept und Lebenslauf hat
die Recherche keine maßgebliche Quelle gefunden.

\subsection{Zu bestätigen}
\begin{enumerate}
  \item Prototype Fund: Geldgeber für J03 (J02: BMFTR, betreut von der Open Knowledge Foundation
    Deutschland; A, SA:9).
  \item Prototype Fund: Beträge. J02-Handbuch: höchstens 950 Stunden und 47.500~EUR für eine
    Person, 1.900 Stunden und 95.000~EUR für Teams, über sechs Monate; Flyer: bis 158.000~EUR in
    zehn Monaten für Teams einschließlich Second Stage (A, SA:11). Second Stage 31.667 und
    63.333~EUR (TV, SA:12, SA:46).
  \item Prototype Fund: Berechtigung. J02: Einzelpersonen oder Teams bis vier, über 18,
    Hauptwohnsitz oder GbR-Sitz in Deutschland; Unternehmen, Hochschulen, Vereine nicht
    berechtigt (A, SA:13). \enquote{In Deutschland steuerpflichtig} steht nicht im Handbuch
    (SA:51).
  \item Prototype Fund: welche Lizenzen anerkannt sind; EUPL-1.2 wird nicht genannt (TV, SA:14).
  \item Prototype Fund: Passung zu Resilienz oder Innovation (FO, Frage 4).
  \item Prototype Fund: Ablauf und Termine nach der Frist für J03; die J02-Daten nicht zitieren
    (TV, SA:16); Formularfragen und Grenzen gegen das J03-Formular (A, SA:6).
  \item § 123 im Land \blank{Land}: zuständige Stelle, Antragsteller, Frist, Kofinanzierung (FO,
    Frage 7). Hessen: Frist (SA:22 nennt keine). Nicht amtlich bestätigt: NRW, Bayern, Baden-Württemberg, Sachsen-Anhalt, Bremen
    (S, SA:24). Niedersachsen: Datum der Richtlinie und Anlagen zum Antragsvordruck (TV, SA:21).
  \item Aktion Mensch: Förderbeträge (nicht recherchiert); Antragsfenster 2027 (S, SA:30); ob die Annahme vorübergehend ausgesetzt ist
    (TV, SA:29); Unterlagen Satzung, Registerauszug, Freistellungsbescheid, Projektbeschreibung,
    Kosten- und Finanzierungsplan und Eigenanteil (S, SA:31).
  \item Deutsche Fernsehlotterie: nur gemeinnützige Träger (S); Freistellungsbescheid und
    Registerauszug für die Organisationsprüfung (S); Stufen von \enquote{Soziale Innovationen}
    (TV); Passung für ein digitales Werkzeug für ältere Zugewanderte (TV) (SA:34).
  \item Vorlage für eine Kooperation oder Weiterleitung mit einem Träger: keine Quelle (TV,
    SA:39).
  \item Civic Coding, Deutsche Stiftung für Engagement und Ehrenamt, Landesprogramme in NRW,
    Berlin und Bayern, Modellvorhaben des BMG, DiPA nach § 40b SGB XI (FO, Fragen 9, 10, 13--15).
\end{enumerate}

% ===========================================================================
\section{Partnerorganisationen}

Quelle: \pfad{scout-organisations.md} (SO, Prüfung vom 23.10.2026, mit Zeilennummer). Rollen:
\textbf{N} nutzt das Werkzeug in der Beratung, \textbf{E} empfiehlt es weiter, \textbf{F} fördert
oder lizenziert möglicherweise. Hier steht nur, was als geprüft gekennzeichnet ist. Das Land des
Betreibers ist offen; örtliche Stellen findet man über die ZQP-Datenbank.

{\small\setlength{\tabcolsep}{3pt}
\begin{longtable}{@{}P{32mm}P{9mm}P{60mm}P{50mm}@{}}
\toprule
\textbf{Name} & \textbf{Rolle} & \textbf{Warum} & \textbf{Zugang}\\
\midrule
\endhead
\multicolumn{4}{@{}l}{\textit{Pflegestützpunkte}}\\
Pflegestützpunkte (örtlich) & N & Beratung nach § 7c SGB XI (FO); kein Geldgeber (SO:13) & über die ZQP-Datenbank \\
ZQP-Beratungsdatenbank & E & Bundesweite Suche nach Stützpunkten und Beratung; berät selbst nicht (SO:8) & \href{https://www.zqp.de/beratung-pflege/}{zqp.de/beratung-pflege} \\
\multicolumn{4}{@{}l}{\textit{Wohlfahrtsverbände}}\\
BAGFW & E, F & Eine Tür zu allen sechs Verbänden (SO:16) & \href{https://www.bagfw.de/site/kontakt}{bagfw.de/site/kontakt} \\
Deutscher Caritasverband & N, E & Pflegeberatung und MBE vor Ort (SO:17) & \href{https://www.caritas.de/kontakt}{caritas.de/kontakt} \\
Diakonie Deutschland & N, E & Pflege und Migration (SO:18) & \href{https://www.diakonie.de/kontakt}{diakonie.de/kontakt} \\
AWO Bundesverband & N, E & Pflege und Migration (SO:19) & \href{https://awo.org/kontakt/}{awo.org/kontakt} \\
DRK & N, E & Pflege, Migration und Teilhabe (SO:20) & \href{https://www.drk.de/das-drk/kontakt/}{drk.de/das-drk/kontakt} \\
Der Paritätische & N, E & Pflege und Migration (SO:21) & \href{https://www.der-paritaetische.de/kontakt/}{der-paritaetische.de/kontakt} \\
ZWST & E & Senioren- und Demenzarbeit; schwächere Passung (SO:22) & \href{https://zwst.org/de/kontakt}{zwst.org/de/kontakt} \\
MBE & N, E & Bester Weg zu türkisch-, russisch-, polnisch- und arabischsprachigen Pflegenden; 1.227 Einrichtungen 2025 (SO:23) & \href{https://www.migrationsberatung.org}{migrationsberatung.org} \\
\multicolumn{4}{@{}l}{\textit{Pflegekassen und Beratung}}\\
GKV-Spitzenverband & E & Spitzenverband der Pflegekassen (SO:26) & \href{https://www.gkv-spitzenverband.de/pflegeversicherung/beratung_und_betreuung/pflegeberatung/pflegeberatung.jsp}{gkv-spitzenverband.de, Pflegeberatung} \\
BARMER & N, F & Eigene Pflegeberatung (SO:28) & \href{https://www.barmer.de/persoenliche-beratung}{barmer.de/persoenliche-beratung} \\
DAK-Gesundheit & \nein & Seite Pflege (SO:30) & \href{https://www.dak.de/dak/kontaktuebersicht_2840}{dak.de, Kontaktübersicht} \\
compass private pflegeberatung & N, F & Gegründet vom PKV-Verband; Telefonberatung für alle Versicherten (SO:31) & \href{https://www.compass-pflegeberatung.de/kontakt}{compass-pflegeberatung.de/kontakt} \\
\multicolumn{4}{@{}l}{\textit{Verbraucher- und Pflegendenverbände}}\\
Verbraucherzentralen, vzbv & E & 16 Landesverbände, etwa 200 Beratungsstellen. Eigener kostenloser Pflegegradrechner mit 64 Kriterien: gleiche Kernfunktion (SO:35) & \href{https://www.verbraucherzentrale.de/ueber-uns}{verbraucherzentrale.de/ueber-uns} \\
BAGSO & E & Etwa 126 Mitgliedsorganisationen (SO:36) & \href{https://www.bagso.de/kontakt/}{bagso.de/kontakt} \\
wir pflegen e.V. & E & Selbsthilfe pflegender Angehöriger, Landesvereine in 10 Ländern (SO:37) & \href{https://www.wir-pflegen.net/kontakt/kontakt-aufnehmen}{wir-pflegen.net, Kontakt} \\
Deutsche Alzheimer Gesellschaft & E & Partnerseite demenz-und-migration.de auf Türkisch, Russisch, Polnisch, Arabisch (SO:38) & \href{https://www.deutsche-alzheimer.de/ueber-uns/kontakt}{deutsche-alzheimer.de, Kontakt} \\
Sozialverband VdK & N, E & Über 2,3 Mio. Mitglieder; Sozialrechtsberatung (SO:40) & \href{https://www.vdk.de/kontakt/}{vdk.de/kontakt} \\
SoVD & N, E & Pflege als Thema (SO:41) & \href{https://www.sovd.de/kontakt}{sovd.de/kontakt} \\
\multicolumn{4}{@{}l}{\textit{Migrantenorganisationen}}\\
Türkische Gemeinde in Deutschland & E & Bietet keine Beratung an; schwache Passung (SO:44) & \href{https://tgd.de/kontakt/}{tgd.de/kontakt} \\
DaMigra & E & Dachverband von Migrantinnenorganisationen (SO:45) & \href{https://www.damigra.de/kontakt/}{damigra.de/kontakt} \\
Verband binationaler Familien & E & Pressemitteilung \enquote{Pflege kennt keine Grenzen} (SO:46) & \href{https://www.verband-binationaler.de/kontakt}{verband-binationaler.de/kontakt} \\
\multicolumn{4}{@{}l}{\textit{Kommunen}}\\
Deutscher Städtetag & E, F & Über die Städte (SO:52) & \href{https://www.staedtetag.de/kontakt}{staedtetag.de/kontakt} \\
Deutscher Landkreistag & E & Keine Kontaktseite gefunden (SO:53) & \href{https://www.landkreistag.de}{landkreistag.de} \\
BaS, Seniorenbüros & E, N & Seniorenbüros bundesweit (SO:54) & \href{https://seniorenbueros.org/kontakt/}{seniorenbueros.org/kontakt} \\
\multicolumn{4}{@{}l}{\textit{Forschung und Qualität}}\\
Stiftung ZQP & E & Mögliche Empfehlung; gegründet vom PKV-Verband (SO:58) & \href{https://www.zqp.de/kontakt/}{zqp.de/kontakt} \\
DZA & \nein & Nur Auswertung oder Forschung (SO:59) & \href{https://www.dza.de/kontakt}{dza.de/kontakt} \\
DIP & \nein & Möglicher bezahlter Auswerter (DIP-GmbH) (SO:61) & \href{https://www.dip.de/ueber-uns/kontakt/}{dip.de, Kontakt} \\
Medizinischer Dienst Bund & \nein & Kein Partner (SO:60) & \href{https://md-bund.de/ueber-uns/ihre-nachricht-an-uns.html}{md-bund.de, Nachricht} \\
\multicolumn{4}{@{}l}{\textit{Digitales und Zivilgesellschaft}}\\
Code for Germany & E & Etwa 20 OK Labs; Hilfe von Freiwilligen, mögliches Zuhause für den Code (SO:64) & \href{https://codefor.de/community/}{codefor.de/community} \\
D64 & \nein & Keine Themen zu Gesundheit oder Pflege; schwache Passung (SO:65) & \href{https://d-64.org/impressum/}{d-64.org} \\
\bottomrule
\end{longtable}
}

\subsection{Was Partner üblicherweise verlangen}
Für einen Pilotversuch mit einem Werkzeug Dritter hat die Recherche bei keiner Gruppe eine
Quelle für eine verlangte Unterlage gefunden (SO:70--73). Mitzubringen sind deshalb, was die
Fragen vorwegnimmt: der Pilotvorschlag, das Factsheet, eine Absichtserklärung zum Unterschreiben,
der Entwurf einer Kooperationsvereinbarung und die Aussagen zu Datenschutz und Barrierefreiheit
(WCAG AAA, durch Tests gesichert; SO:73). Jeder Partner wird fragen, ob die vier Fremdsprachen
von Muttersprachlern geprüft sind; sie sind es nicht (SO:5). Die Verbraucherzentralen werden
fragen, warum nicht ihr eigener Rechner; Anspruch unterscheidet sich durch sechs Sprachen, die
Berechnung der Lücke, das PDF und die kurze Befragung (SO:35). Eine Pflegekasse als Geldgeber
hätte nach Einschätzung der Recherche einen Interessenkonflikt, weil das Werkzeug Familien zu
höheren Ansprüchen verhilft; dafür gibt es keine Quelle (SO:32).

\subsection{Empfohlene erste drei für einen Pilotversuch}
\begin{enumerate}
  \item \textbf{Ein Pflegestützpunkt in \blank{Land}}, gefunden über die ZQP-Datenbank: berät
    nach § 7c SGB XI genau die Familien, die das Werkzeug zur Beratung schickt (SO:8, SO:13; FO).
  \item \textbf{Eine MBE-Stelle eines Wohlfahrtsverbands} (Caritas, Diakonie, AWO, DRK oder
    Paritätischer) in \blank{Land}: der beste Weg zu Pflegenden, die Türkisch, Russisch, Polnisch
    oder Arabisch sprechen (SO:23).
  \item \textbf{Eine Beratungsstelle des Sozialverbands VdK}: berät selbst im Sozialrecht und
    kann das Werkzeug in der Beratung nutzen und weiterempfehlen (SO:40).
\end{enumerate}

\subsection{Zu bestätigen}
\begin{enumerate}
  \item Ältere Adresse beratungsdatenbank.zqp.de; ob die ZQP-Datenbank alle Pflegestützpunkte
    führt (SO:8).
  \item Rahmenverträge nach § 7c zwischen Landesverbänden der Pflegekassen und Ländern (SO:9).
  \item Pflegeberatungs-Richtlinien nach § 17 Abs.~1a SGB XI und ob ein Werkzeug Dritter mit
    ihnen vereinbar ist; Endfassung nicht gefunden (SO:10, SO:26).
  \item Landesarbeitsgemeinschaften der Pflegestützpunkte (Rheinland-Pfalz,
    Baden-Württemberg), Adressliste Bayern, Netzwerke in Sachsen und Sachsen-Anhalt (SO:11--12).
  \item KI-Pilotwerkstatt des Paritätischen und ihre Datenschutzanforderungen (SO:21, SO:71).
  \item AOK-Bundesverband: Seite und Kontakt (SO:27); TK: Kontaktseite (SO:29).
  \item BIVA-Pflegeschutzbund: Kontakt und Angebote (SO:39).
  \item BV NeMO, etwa 800 Organisationen in 21 Netzwerken (SO:47); Forum für eine
    kultursensible Altenhilfe, möglicherweise ruhend (SO:48).
  \item Bündnis Digitalisierung in der Pflege: Kontakt (SO:66).
  \item Widerspruchsfunktion des Pflegegradrechners der Verbraucherzentralen seit August 2025
    (SO:35); MBE-Standortsuche (BAMF NAvI) (SO:23).
\end{enumerate}

% ===========================================================================
\section{Unterlagen}

\subsection{Matrix}
Spalten: PF = Prototype Fund; 123 = § 123-Vorhaben über einen Träger; AM = Aktion Mensch über
einen Träger; Pilot = Pilotversuch mit einem Partner; Recht = Betrieb der Seite.
\ja{} gebraucht, \nuetzlich{} nützlich, \nein{} nicht nötig. Status: \textbf{da} = im Repository,
\textbf{Entwurf} = in diesem Lauf entworfen, \textbf{Lücken} = mit auszufüllenden Stellen,
\textbf{fehlt}. Pfade von Entwürfen relativ zu \pfad{docs/dossier/}.

{\footnotesize\setlength{\tabcolsep}{3pt}
\begin{longtable}{@{}P{31mm}*{5}{>{\centering\arraybackslash}p{9mm}}P{19mm}P{46mm}@{}}
\toprule
\textbf{Unterlage} & \textbf{PF} & \textbf{123} & \textbf{AM} & \textbf{Pilot} & \textbf{Recht}
  & \textbf{Status} & \textbf{Pfad}\\
\midrule
\endhead
Impressum (Seite) & \nuetzlich & \nuetzlich & \nuetzlich & \ja & \ja & da, Lücken &
  \pfad{app/impressum/}, \pfad{lib/legal/impressum.ts} \\
Datenschutzerklärung (Seite) & \nuetzlich & \nuetzlich & \nuetzlich & \ja & \ja & da, Lücken &
  \pfad{app/datenschutz/}, \pfad{lib/legal/datenschutz.ts} \\
Compliance-Liste & \nein & \nuetzlich & \nein & \nuetzlich & \ja & da &
  \pfad{docs/COMPLIANCE.md} \\
Pilotvorschlag & \nein & \nuetzlich & \nein & \ja & \nein & da, Lücken &
  \pfad{docs/proposal-organisations.md} \\
Förderübersicht & \ja & \ja & \ja & \nein & \nein & da, Lücken & \pfad{docs/funding-outline.md} \\
Projektbeschreibung & \ja & \ja & \ja & \nuetzlich & \nein & Entwurf, Lücken &
  \pfad{antrag/projektbeschreibung.tex} \\
Kosten- und Finanzierungsplan & \nuetzlich & \ja & \ja & \nein & \nein & Entwurf, Lücken &
  \pfad{antrag/finanzplan.tex} \\
Zeit- und Meilensteinplan & \ja & \ja & \ja & \nuetzlich & \nein & Entwurf, Lücken &
  \pfad{antrag/zeitplan.tex} \\
Anschreiben & \nein & \ja & \ja & \ja & \nein & Entwurf, Lücken &
  \pfad{partner/anschreiben.tex} \\
Letter of Intent & \nuetzlich & \ja & \ja & \ja & \nein & Entwurf, Lücken &
  \pfad{partner/letter-of-intent.tex} \\
Kooperationsvereinbarung & \nein & \ja & \ja & \ja & \nuetzlich & Entwurf, Lücken &
  \pfad{partner/kooperationsvereinbarung-pilot.tex} \\
Factsheet & \nuetzlich & \nuetzlich & \nuetzlich & \ja & \nein & Entwurf, Lücken &
  \pfad{partner/factsheet.tex} \\
Datenschutzkonzept & \nein & \nuetzlich & \nuetzlich & \ja & \nuetzlich & fehlt & \nein \\
Barrierefreiheitserklärung & \nein & \nuetzlich & \nuetzlich & \nuetzlich & \ja & fehlt & \nein \\
Lebenslauf (Gründer) & \nuetzlich & \nuetzlich & \nein & \nein & \nein & fehlt & \nein \\
Freistellungsbescheid & \nein & \nein & \ja & \nein & \nein & nur über Träger & \nein \\
Lizenzdatei & \ja & \nuetzlich & \nein & \nuetzlich & \ja & fehlt & \texttt{LICENSE} (geplant) \\
AVV Vercel & \nein & \nein & \nein & \nuetzlich & \ja & abzuschließen & \nein \\
AVV Plausible & \nein & \nein & \nein & \nuetzlich & \ja & abzuschließen & \nein \\
VVT (Art.~30 DSGVO) & \nein & \nein & \nein & \nuetzlich & \ja & fehlt & \nein \\
Dieses Dossier & \nuetzlich & \nuetzlich & \nuetzlich & \nuetzlich & \nein & Entwurf, Lücken &
  \pfad{anspruch-dossier.tex} \\
\bottomrule
\end{longtable}
}

Grundlage der Zuordnung: Prototype Fund ohne Anhänge, nur Text und Links, mit
Open-Source-Zusage (SA:10, SA:14); § 123 mit Konzeption, Finanzierungsplan und Antrag in
Textform (SA:19); Aktion Mensch mit Bescheid des Finanzamts (SA:29), weitere Unterlagen zu
bestätigen (SA:31); Stiftungsleitfaden (SA:43); Recht nach \pfad{docs/COMPLIANCE.md}. Für den
Pilotversuch gibt es keine Quelle; die Zuordnung folgt dem Pilotvorschlag. Ein
\nuetzlich{} ist eine Einschätzung, keine Anforderung eines Geldgebers.

\subsection{Was die fehlenden Unterlagen enthalten sollen}
\begin{description}
  \item[Datenschutzkonzept.] Für Partner und Antrag, auf zwei bis drei Seiten: welche Daten wo
    entstehen (Antworten nur im Browser; Hosting-Protokolle bei Vercel; Zählung bei Plausible,
    falls an; E-Mail-Kontakt); was nie erhoben wird; Rechtsgrundlagen wie in der
    Datenschutzerklärung; Auftragsverarbeiter und Übermittlungen; Speicherdauer; technische
    Maßnahmen (Erlaubnisliste, Tests, kein Cookie, Partnerkürzel nur im Speicher); Rechte der
    Betroffenen und Widerspruch; Spracheingabe; Rollen im Pilotversuch (keine Übermittlung
    personenbezogener Daten zwischen Partner und Betreiber, zu bestätigen). Eine maßgebliche
    Vorlage hat die Recherche nicht gefunden (SA:43).
  \item[Barrierefreiheitserklärung.] Für öffentliche Stellen, die die Seite einbinden und unter
    BITV 2.0 fallen, auf Anfrage (COMPLIANCE 2.9). Grundlage sind die WCAG-AAA-Kontraste und
    Tests aus \pfad{README.md} und die bekannten Grenzen (Spracheingabe nur in manchen
    Browsern, lange Vorlesezeit). Die Pflichtbestandteile nach BITV 2.0 sind zu bestätigen.
  \item[Lebenslauf.] Für den Gründer, eine Seite: Ausbildung, Erfahrung in Software, bis drei
    Projekte (im J02-Formular als \enquote{Erfahrung} gefragt; A, SA:10).
  \item[Freistellungsbescheid.] Kann der Gründer nicht vorlegen; nur ein gemeinnütziger Träger hat
    ihn (SA:29).
  \item[Lizenzdatei.] \texttt{LICENSE} mit dem Text der EUPL-1.2, wenn die Empfehlung angenommen
    wird; dazu eine Vereinbarung für Beiträge Dritter (FO).
  \item[AVV.] Vercel: prüfen, ob der Vertrag den Plan abdeckt; Plausible: durch Nutzung
    angenommen, ablegen (COMPLIANCE 2.4, 2.5).
  \item[VVT.] Eine Seite, drei Einträge: Vercel-Protokolle, Plausible-Statistik, E-Mail-Kontakt
    aus dem Impressum (COMPLIANCE 2.8).
\end{description}

% ===========================================================================
\section{Die nächsten 90 Tage}

Woche 1 beginnt am 23.10.2026. Die einzige feste Frist ist der 30.11.2026 (Prototype Fund,
Woche 6). Termine Dritter stehen zur Orientierung dabei.

{\small\setlength{\tabcolsep}{3pt}
\begin{longtable}{@{}P{22mm}P{110mm}P{20mm}@{}}
\toprule
\textbf{Woche} & \textbf{Schritt} & \textbf{Grundlage}\\
\midrule
\endhead
1 (bis 29.10.) & Lizenz entscheiden (Empfehlung EUPL-1.2). Fragen an den Prototype Fund zu
  Berechtigung, Passung und Lizenz senden. Land festlegen und Aufsichtsbehörde bestimmen. & FO;
  COMPL. 2.1 \\
1--2 & Lücken in \pfad{lib/legal/blanks.ts} füllen. Muttersprachler für die vier Sprachen
  suchen und beauftragen. Kurze rechtliche Prüfung anfragen (DSK Rn.~87--90, 68/69, 80). &
  COMPL. 1, 2.6, 2.7; LC:32 \\
2 (bis 05.11.) & Vercel-Plan und Plausible-Tarif wählen; beide AVV prüfen und ablegen. Die drei
  Fragen zum Quelltext entscheiden. PR \#7 und \#8 zusammenführen. & COMPL. 2.2--2.5; LC:34 \\
2--3 & Erste Lesung PNOG am 05.11.2026 verfolgen (Termin Dritter). & LC:26 \\
3 & Impressum und Datenschutzerklärung mit gefüllten Lücken live. \texttt{LICENSE} einchecken.
  VVT schreiben. Partnerlink auf der Produktionsadresse prüfen. & COMPL. 3.11, 2.8;
  \pfad{DEPLOY.md} \\
3--5 & Bewerbung Prototype Fund schreiben, wenn die Antworten es tragen, aus
  \pfad{antrag/projektbeschreibung.tex} und \pfad{antrag/zeitplan.tex}. & SA:10 \\
4 & Preis des Pilotversuchs festlegen (\blank{Betrag} oder kostenlos). Die ersten drei Partner
  anschreiben (Anschreiben, Factsheet, Pilotvorschlag). & SO; UP:30 \\
5 (bis 26.11.) & Bewerbung Prototype Fund einreichen, mit Puffer vor der Frist. & SA:7 \\
6 (30.11.) & \textbf{Frist Prototype Fund.} & SA:7 \\
5--8 & Gespräche mit Partnern; Absichtserklärungen einholen; Kooperationsvereinbarung zur
  anwaltlichen Prüfung. Zählung einschalten, wenn COMPLIANCE Abschnitt 3 erfüllt ist. &
  COMPL. 3 \\
6--8 & § 123: Landesministerium und Landesverbände der Pflegekassen in \blank{Land} fragen;
  Gespräch mit einem möglichen Träger (\blank{Kommune oder Träger}). Fernsehlotterie: Bewerbung
  bis 17.12.2026 nur, wenn ein Träger und die Passung feststehen (Termin Dritter). & FO Frage
  7; SA:34 \\
8--10 & Falls das PNOG beschlossen wird: Beträge und Schwellen ab 01.01.2027 anpassen.
  Prüfung durch Muttersprachler abschließen. & LC:26, 36 \\
11 (ab 01.01.) & Aktion Mensch: neue Antragsphase prüfen, nur mit Träger (Termin Dritter). & SA:27 \\
11--13 (bis 21.01.) & Erster Pilotversuch beginnt am \blank{Startdatum}: zwei Beratende, ein
  Gespräch im Monat, Strichliste ohne Namen. & Pilot\-vorschlag \\
\bottomrule
\end{longtable}
}

% ===========================================================================
\section{Anhang: Quellen}

\subsection*{Im Repository}
{\small\raggedright
\pfad{README.md}; \pfad{docs/COMPLIANCE.md}; \pfad{docs/DEPLOY.md};
\pfad{docs/PARTNERS.md}; \pfad{docs/funding-outline.md};
\pfad{docs/proposal-organisations.md}; \pfad{lib/rules/sources.ts};
\pfad{lib/legal/datenschutz.ts}; \pfad{lib/legal/blanks.ts}; die Zusammenfassungen in
\nolinkurl{.orchestration/runs/2026-10-22-*/99_SUMMARY.md}; die Recherchen
\pfad{scout-applications.md} und \pfad{scout-organisations.md} in
\pfad{.orchestration/runs/2026-10-23-dossier/}.\par}

\subsection*{Förderprogramme}
{\scriptsize\raggedright
\textbullet~\url{https://bewerben.prototypefund.de/}
\textbullet~\url{https://bewerben.prototypefund.de/jahrgang-03-resilienz/}
\textbullet~\url{https://www.prototypefund.de/uploads/Publikationen_und_Onepager/HandbuchBewerbung.pdf} (archiviert)
\textbullet~\url{https://kb.prototypefund.de/books/rechtliche-grundlagen}
\textbullet~\url{https://www.gesetze-im-internet.de/sgb_11/__123.html}
\textbullet~\url{https://www.gkv-spitzenverband.de/media/dokumente/pflegeversicherung/richtlinien__vereinbarungen__formulare/2026-06-01_Pflege_Empfehlungen_nach_123_Abs.3_SGB_XI.pdf}
\textbullet~\url{https://www.berlin.de/sen/pflege/grundlagen/modellvorhaben-nach----123--124-sgb-xi-im-land-berlin-1687308.php}
\textbullet~\url{https://soziales.niedersachsen.de/gemeinsame-modellvorhaben/richtlinie-uber-die-gewahrung-von-zuwendungen-zur-forderung-gemeinsamer-modellvorhaben-nach-123-sgb-xi-245717.html}
\textbullet~\url{https://familie.hessen.de/sites/familie.hessen.de/files/2025-10/erlaeuterungen_zum_antragsformular-modellprojekte_nach_ss_123_sgb_xi.pdf}
\textbullet~\url{https://www.schleswig-holstein.de/DE/fachinhalte/P/pflege/_documents/modellvorhaben_pflegebeduerftige}
\textbullet~\url{https://www.aok.de/pp/nordwest/pm/foerderbescheid-projekt-zapfdigital-sozialministerin-schleswig-holstein-2026/}
\textbullet~\url{https://www.land.nrw/pressemitteilung/foerderung-fuer-modellvorhaben-zur-unterstuetzung-und-stabilisierung-haeuslicher}
\textbullet~\url{https://www.aktion-mensch.de/foerderung/foerderprogramme}
\textbullet~\url{https://www.aktion-mensch.de/foerderung/antrag/foerderfinder}
\textbullet~\url{https://www.aktion-mensch.de/foerderung/angebote/schnell-check}
\textbullet~\url{https://zwst.org/sites/default/files/2026-08/AktionMensch_ZWST_5_6_2026_\%20end.pdf}
\textbullet~\url{https://www.fernsehlotterie.de/foerdern-engagieren/foerdermittelbewerbung}
\textbullet~\url{https://www.fernsehlotterie.de/neuigkeiten-stiftung/mehrstufiges-bewerbungsverfahren}
\textbullet~\url{https://foerderportal.deutsches-hilfswerk.de/}
\textbullet~\url{https://www.bay-landkreistag.de/media/103712/vwi-06092024-informationsveranstaltung-stiftung-dt-hilfswerk.pdf}
\textbullet~\url{http://bosch-stiftung.de/de/foerderung}
\textbullet~\url{https://www.bosch-stiftung.de/en/presse/2026}
\textbullet~\url{https://www.stiftungen.org/fileadmin/stiftungen_org/Verband/Was_wir_tun/Veranstaltungen/AK-Foerderstiftungen/Leitfaden-fuer-Foerderanfragen-an-Stiftungen.pdf}
\textbullet~\url{https://www.civic-innovation.de/fileadmin/Dokumente/CIP_Foerderrichtlinie_2025_2026.pdf}
\textbullet~\url{https://www.der-paritaetische.de/alle-meldungen/neue-foerderkriterien-des-deutschen-hilfswerks-fuer-investitionen-in-baumassnahmen-undoder-erstausstat/}
\textbullet~\url{https://grantedai.com/grants/prototype-fund-open-knowledge-foundation-germany-ac23a4d2} (sekundär)
\textbullet~\url{https://interoperable-europe.ec.europa.eu/collection/eupl/faq}
\par}

\subsection*{Partnerorganisationen}
{\scriptsize\raggedright
\textbullet~\url{https://www.zqp.de/beratung-pflege/}, \url{https://www.zqp.de/kontakt/},
        \url{https://www.zqp.de/ueber-uns/kooperation/}
\textbullet~\url{https://lfp.bayern.de}
\textbullet~\url{https://www.bagfw.de}, \url{https://www.bagfw.de/site/kontakt}
\textbullet~\url{https://www.caritas.de/kontakt},
        \url{https://www.caritas.de/hilfeundberatung/ratgeber/alter/pflege/pflege},
        \url{https://www.caritas.de/fuerprofis/fachthemen/migration/fachthema-migration}
\textbullet~\url{https://www.diakonie.de/kontakt},
        \url{https://www.diakonie.de/informieren/unsere-themen/gesundheit-pflege},
        \url{https://www.diakonie.de/informieren/unsere-themen/flucht-migration}
\textbullet~\url{https://awo.org/kontakt/}, \url{https://awo.org/themen/gesundheit-pflege-und-teilhabe/},
        \url{https://awo.org/themen/migration-und-flucht/}
\textbullet~\url{https://www.drk.de/das-drk/kontakt/},
        \url{https://www.drk.de/hilfe-in-deutschland/senioren/pflege-und-betreuung/},
        \url{https://www.drk.de/hilfe-in-deutschland/migration-integration-und-teilhabe/}
\textbullet~\url{https://www.der-paritaetische.de/kontakt/}
\textbullet~\url{https://zwst.org/de/kontakt}
\textbullet~\url{https://www.migrationsberatung.org}
\textbullet~\url{https://www.gkv-spitzenverband.de/pflegeversicherung/beratung_und_betreuung/pflegeberatung/pflegeberatung.jsp}
\textbullet~\url{https://www.barmer.de/persoenliche-beratung}
\textbullet~\url{https://www.tk.de}
\textbullet~\url{https://www.dak.de/dak/kontaktuebersicht_2840}
\textbullet~\url{https://www.compass-pflegeberatung.de/kontakt}
\textbullet~\url{https://www.verbraucherzentrale.de/ueber-uns},
        \url{https://www.verbraucherzentrale.de/pflegegradrechner}
\textbullet~\url{https://www.bagso.de/kontakt/}
\textbullet~\url{https://www.wir-pflegen.net/kontakt/kontakt-aufnehmen}
\textbullet~\url{https://www.deutsche-alzheimer.de/ueber-uns/kontakt}
\textbullet~\url{https://www.biva.de}
\textbullet~\url{https://www.vdk.de/kontakt/}
\textbullet~\url{https://www.sovd.de/kontakt}
\textbullet~\url{https://tgd.de/kontakt/}
\textbullet~\url{https://www.damigra.de/kontakt/}
\textbullet~\url{https://www.verband-binationaler.de/kontakt}
\textbullet~\url{https://www.staedtetag.de/kontakt}
\textbullet~\url{https://www.landkreistag.de}
\textbullet~\url{https://seniorenbueros.org/kontakt/}
\textbullet~\url{https://www.dza.de/kontakt}
\textbullet~\url{https://md-bund.de/ueber-uns/ihre-nachricht-an-uns.html}
\textbullet~\url{https://www.dip.de/ueber-uns/kontakt/}
\textbullet~\url{https://codefor.de/community/}
\textbullet~\url{https://d-64.org/impressum/}
\par}

\subsection*{Recht, Datenschutz und Betrieb}
{\scriptsize\raggedright
\textbullet~\url{https://www.gesetze-im-internet.de/sgb_11/__7a.html},
        \url{https://www.gesetze-im-internet.de/sgb_11/__7b.html},
        \url{https://www.gesetze-im-internet.de/sgb_11/__7c.html},
        \url{https://www.gesetze-im-internet.de/sgb_11/__45.html},
        \url{https://www.gesetze-im-internet.de/sgb_11/__40b.html}
\textbullet~\url{https://www.bundesgesundheitsministerium.de/fileadmin/Dateien/3_Downloads/Gesetze_und_Verordnungen/GuV/P/Kabinett-GE_Pflegeneuordnungsgesetz_PNOG.pdf}
\textbullet~\url{https://www.aok.de/gp/pflegebeduerftigkeit/besondere-bedarfskonstellation}
\textbullet~\url{https://www.bfdi.bund.de/DE/Service/Anschriften/Laender/Laender-node.html}
\textbullet~\url{https://datenschutzkonferenz-online.de/media/oh/OH_Digitale_Dienste.pdf}
\textbullet~\url{https://vercel.com/legal/privacy-policy}, \url{https://vercel.com/legal/dpa}
\textbullet~\url{https://vercel.com/docs/analytics/custom-events},
        \url{https://vercel.com/docs/analytics/limits-and-pricing}
\textbullet~\url{https://plausible.io/docs/custom-props/introduction},
        \url{https://plausible.io/docs/shared-links}
\par}

\end{document}
